% Generated draft: requires mathematical and source review.
\documentclass[11pt]{article}
\usepackage[margin=25mm]{geometry}
\usepackage{fontspec}
\setmainfont{Latin Modern Roman}
\usepackage{amsmath,amssymb,mathtools,unicode-math}
\setmathfont{Latin Modern Math}
\usepackage{enumitem,needspace,xurl,hyperref}
\setlength{\parindent}{0pt}
\setlength{\parskip}{5pt}
\setlength{\emergencystretch}{3em}
\newcommand{\N}{\mathbb{N}}
\newcommand{\Z}{\mathbb{Z}}
\newcommand{\Q}{\mathbb{Q}}
\newcommand{\R}{\mathbb{R}}
\newcommand{\C}{\mathbb{C}}
\newcommand{\F}{\mathbb{F}}
\newcommand{\A}{\mathbb{A}}
\newcommand{\PP}{\mathbb{P}}
\newcommand{\E}{\mathbb{E}}
\newcommand{\eps}{\varepsilon}
\newcommand{\dd}{\,\mathrm{d}}
\begin{document}
\section*{UnsolvedMath ID: 2800406}
\textbf{Source problem code:} catalogue entry
\subsection{Definitions and mathematical statement}
Let \(n\ge1\) be an integer, and let \(X_1,\ldots,X_n\) be independent random variables on a probability space with
\[
X_i\ge0\text{ almost surely},\qquad \mathbb E X_i=1\quad(1\le i\le n).
\]
No identical-distribution assumption is imposed. Feige's conjecture as recorded here asserts
\[
\Pr\left\{\sum_{i=1}^nX_i\ge n+1\right\}\le1-e^{-1}
\]
for every such \(n\) and collection. The threshold includes equality.

The supplied research summary reports a July 2026 resolution with the sharper, finite-\(n\) bound
\[
\Pr\left\{\sum_{i=1}^nX_i\ge n+1\right\}
\le1-\left(\frac n{n+1}\right)^n\le1-e^{-1},
\]
which it calls sharp. It describes Fu--Han--Wang--Yan--Zhang--Zhou as covering a general-shift regime \(\delta\ge1\), and Nie--Wei as covering \(\delta\le1\), with both relevant at the stated unit shift. No further general-shift formula is supplied or asserted here.

This is explicitly a supplied preprint assessment: the summary calls the papers about one week old at the time of its assessment, not peer-reviewed, and AI-assisted with disclosure. It reports two independent proofs and a Lean formalization as grounds for its high confidence. None of those proof, independence, formalization, or review claims has been independently checked in this draft.
\subsection{Short English statement}
Independent nonnegative mean-one random variables should exceed their total mean by at least one with probability at most \(1-e^{-1}\). The supplied summary reports a sharper finite-dimensional resolution in recent, unreviewed July 2026 preprints.
\subsection{Sources}
\begin{itemize}
\item Local UnsolvedMath registry, permanent ID \(2800406\), exact source code \texttt{catalogue entry}, ``10 Lectures and 42 Open Problems --- Feige's conjecture.'' Proposed-by metadata: Afonso S. Bandeira, 2015; conjecture attributed to Feige.
\item Bandeira, \emph{42 Open Problems in Data Science} (2016), Open Problem 4.6; author-blog source \url{https://afonsobandeira.wordpress.com/2015/11/29/10l42panextraopenproblem/}. HTML extraction; recorded access 2026-07-29.
\item Supplied research summary: Fu--Han--Wang--Yan--Zhang--Zhou, arXiv:2607.23980, reported sharp result and Lean verification; Nie--Wei, arXiv:2607.24528, reported independent result. The summary states that these build on Vlassis--Thomas, arXiv:2607.08415 (Gaffke's conjecture), and Letwin--Yaskin, arXiv:2410.04741. No independent preprint or formalization audit is claimed.
\end{itemize}
\end{document}
